\documentclass[11pt]{article}

% --- Packages ---
\usepackage[margin=1in]{geometry}
\usepackage{amsmath,amssymb}
\usepackage{graphicx}
\usepackage[hidelinks]{hyperref}
\usepackage{natbib}
\usepackage{booktabs}
\usepackage{enumitem}
\usepackage{xcolor}
\usepackage{siunitx}
\usepackage{subcaption}
\usepackage{titlesec}
\usepackage{setspace}

% --- Formatting ---
\titleformat{\section}{\normalfont\Large\bfseries}{\thesection}{1em}{}
\titleformat{\subsection}{\normalfont\large\bfseries}{\thesubsection}{1em}{}
\setlength{\parskip}{0.5em}
\setlength{\parindent}{0em}

% --- Draft helpers (remove for final) ---
\newcommand{\todo}[1]{\textcolor{red}{\textbf{[TODO: #1]}}}
\newcommand{\note}[1]{\textcolor{blue}{\textbf{[NOTE: #1]}}}

% =============================================================================
\begin{document}

\begin{center}
    {\LARGE \textbf{Bayesian Optimization of Tensegrity Structures \\
    for Energy Absorption: A Simulation-Guided \\
    Experimental Approach}} \\[1em]
    {\large BYU College of Engineering --- Mentored Research Grant Proposal} \\[0.5em]
    {\large PI: \todo{Faculty Name}, Department of \todo{Department}} \\
    {\large Duration: 2 Years \quad | \quad Requested: \$25,000}
\end{center}

\vspace{1em}
\hrule
\vspace{1em}

% =============================================================================
\section{Project Summary}
% =============================================================================

We propose a two-year mentored research program in which 2--3 undergraduate
students investigate the design of tensegrity structures optimized for energy
absorption using a closed-loop framework that integrates physics-based
simulations with Bayesian optimization and physical validation experiments.

Tensegrity (tensional integrity) structures---assemblies of rigid bars held in
stable equilibrium by a network of tensioned cables---exhibit remarkable
strength-to-weight ratios and tunable mechanical responses.  By varying
structural parameters such as bar lengths, cable stiffnesses, angles, and
induced curvature, a rich and high-dimensional design space emerges.  Navigating
this space efficiently requires intelligent experimental design; Bayesian
optimization provides exactly this capability by constructing a surrogate model
of the structure--property relationship and guiding subsequent experiments toward
optimal configurations.

The project emphasizes \textbf{undergraduate mentoring} as its primary
objective.  Students will gain hands-on experience in experimental mechanics,
data acquisition, physics-based simulation, statistical modeling, CAD design, and
scientific communication---skills directly transferable to graduate school and
industry careers.

% =============================================================================
\section{Research Objectives}
% =============================================================================

\begin{enumerate}[label=\textbf{O\arabic*.}]
    \item \textbf{Define and parameterize a tensegrity unit cell} with
        systematically variable design parameters (bar length, bar thickness,
        cable length, cable stiffness, number of bars/cables, bar/cable angles,
        and curvature/induced buckling).

    \item \textbf{Develop physics-based simulation models} (finite element or
        analytical) of the tensegrity unit cell that predict force transmission,
        deformation, and energy absorption under impact loading.

    \item \textbf{Implement a Bayesian optimization loop} that uses simulation
        predictions as the objective function to efficiently explore the design
        space, identifying configurations that maximize energy absorption and
        minimize transmitted force.

    \item \textbf{Fabricate and experimentally validate} a set of approximately
        five tensegrity structures, comparing measured force--time histories and
        deformation responses against simulation predictions (ex-situ
        validation).

    \item \textbf{Characterize energy absorption performance} through
        controlled drop tests at increasing heights, identifying the maximum
        drop height sustainable without plastic deformation.

    \item \textbf{(Stretch)} Quantify advanced metrics including specific energy
        absorption (\si{J/kg}), effective spring constant per unit weight, and
        effective spring constant per unit volume.
\end{enumerate}

% =============================================================================
\section{Background and Motivation}
% =============================================================================

\subsection{Tensegrity Structures}

Tensegrity structures, first described by Buckminster Fuller and artist Kenneth
Snelson, consist of isolated compression members (bars or struts) suspended
within a continuous network of tension members (cables or strings)
\citep{skelton2009tensegrity, sultan2009tensegrity}.  Their defining
characteristic---that no two rigid bars touch---yields lightweight assemblies
with surprising load-bearing capacity and resilience.

The mechanical response of a tensegrity structure is governed by a large number
of design variables:
\begin{itemize}
    \item \textbf{Bar parameters:} length, cross-sectional thickness, material
        stiffness, curvature (enabling induced buckling modes).
    \item \textbf{Cable parameters:} length, cross-sectional thickness,
        material stiffness (elastic modulus, yield strain).
    \item \textbf{Topology:} number of bars and cables, connectivity pattern.
    \item \textbf{Geometry:} angles between bars and cables, overall aspect
        ratio, boundary conditions.
\end{itemize}

This results in a combinatorially large design space, making exhaustive
experimental search impractical and motivating data-efficient optimization
strategies.

\subsection{Bayesian Optimization for Structural Design}

Bayesian optimization (BO) is a sequential, model-based strategy for optimizing
expensive-to-evaluate black-box functions \citep{shahriari2016taking,
frazier2018tutorial}.  A Gaussian process (GP) surrogate model is fitted to
observed data and an acquisition function (e.g., Expected Improvement) balances
exploration of uncertain regions with exploitation of promising areas.  BO has
been successfully applied to materials discovery \citep{lookman2019active},
mechanical metamaterial design \citep{wang2022bayesian}, and structural
optimization \citep{lee2023bayesian}.

Incorporating physics-based simulations into the BO loop---sometimes called
simulation-in-the-loop optimization---enables rapid surrogate updates without
requiring a physical experiment for every candidate design
\citep{perdikaris2017nonlinear}.  Physical experiments are reserved for
validation, dramatically reducing cost and time.

\subsection{Energy Absorption and Impact Mitigation}

Quantifying energy absorption under impact is central to protective structure
design.  Key metrics include:
\begin{itemize}
    \item \textbf{Peak transmitted force:} the maximum force recorded on the
        distal side of the structure when subjected to an impact (e.g., a ball
        drop).
    \item \textbf{Force--time profile:} the full temporal evolution of
        transmitted force, from which impulse and energy dissipation can be
        computed.
    \item \textbf{Deformation:} total displacement and residual (plastic)
        deformation after impact.
    \item \textbf{Critical drop height:} the maximum height from which the
        structure can be impacted without exhibiting plastic (permanent)
        deformation, determined via progressive drop tests starting from low
        heights.
\end{itemize}

% =============================================================================
\section{Technical Approach}
% =============================================================================

\subsection{Phase~1: Design Space Definition and Simulation (Months 1--8)}

\subsubsection{Tensegrity Unit Cell Parameterization}

We will define a parameterized tensegrity unit cell suitable for both simulation
and fabrication.  The baseline topology will be a three-bar prismatic tensegrity
\citep{skelton2009tensegrity}, chosen for its relative simplicity and
well-understood mechanics, while still offering sufficient design freedom.

The parameterization will include:
\begin{itemize}
    \item Bar lengths ($L_b$), thicknesses ($t_b$), and material properties
    \item Cable lengths ($L_c$), thicknesses ($t_c$), and stiffnesses ($k_c$)
    \item Number of bars ($n_b$) and cables ($n_c$)
    \item Bar twist angle ($\alpha$) and cable attachment angles ($\beta$)
    \item Induced bar curvature ($\kappa$) to explore pre-buckled
        configurations
\end{itemize}

Each parameter will be assigned physically meaningful bounds based on
manufacturing constraints and material availability.

\subsubsection{Physics-Based Simulation}

We will develop finite element models of the tensegrity unit cell using
open-source tools (e.g., FEniCS, or a combination of Python with analytical
stiffness matrix methods).  The simulation will:
\begin{enumerate}
    \item Accept the parameterized design vector as input.
    \item Apply a transient impact load representative of a ball-drop
        experiment.
    \item Output the force--time history at the base, peak transmitted force,
        total deformation, and absorbed energy.
\end{enumerate}

Simulation fidelity will be validated iteratively against early experimental
results (see Phase~2).

\subsubsection{Bayesian Optimization Implementation}

A Bayesian optimization loop will be constructed using established Python
libraries (e.g., BoTorch, Ax, or a custom GP implementation).  The simulation
model serves as the objective function:

\begin{equation}
    \mathbf{x}^* = \arg\min_{\mathbf{x} \in \mathcal{X}} f_{\text{sim}}(\mathbf{x})
\end{equation}

\noindent where $\mathbf{x}$ is the design parameter vector, $\mathcal{X}$ is
the feasible design space, and $f_{\text{sim}}(\mathbf{x})$ returns the peak
transmitted force (or a composite objective incorporating deformation and energy
metrics).  The optimization will proceed for a predetermined budget of
simulation evaluations, after which the top-performing candidates will be
selected for fabrication.

\subsection{Phase~2: Fabrication and Experimental Validation (Months 6--20)}

\subsubsection{Tensegrity Fabrication}

Approximately five tensegrity structures will be fabricated based on the
Bayesian-optimized designs and selected comparison baselines.  Structures will
be:
\begin{itemize}
    \item Designed in CAD software (e.g., SolidWorks or Fusion 360) to ensure
        dimensional accuracy and reproducibility.
    \item Fabricated using a combination of 3D printing (for rigid bar
        components) and commercially available cables/strings of known
        stiffness.
    \item Assembled with careful attention to dimensional tolerances, as
        accuracy is critical for meaningful comparison with simulations.
\end{itemize}

\subsubsection{Impact Testing Protocol}

Each structure will undergo a controlled impact test:
\begin{enumerate}
    \item A ball of known mass is dropped from a controlled height onto the top
        of the tensegrity structure.
    \item A force sensor (load cell) at the base records the transmitted
        force--time history.
    \item High-speed video (if available) or displacement sensors capture
        deformation.
    \item Tests begin at low drop heights and increase incrementally until
        plastic deformation is detected, establishing the critical drop height.
\end{enumerate}

\textbf{Experimental rigor:} Students will be trained in proper experimental
methodology, including:
\begin{itemize}
    \item Background and baseline measurements (e.g., force sensor calibration,
        drop without a structure).
    \item Repeatability studies (multiple drops per configuration at the same
        height).
    \item Systematic documentation of experimental conditions and anomalies.
\end{itemize}

\subsubsection{Ex-Situ Validation}

Experimental results will be compared quantitatively against simulation
predictions for the same design parameters.  Discrepancies will inform
simulation model refinement and uncertainty quantification in the GP surrogate.

\subsection{Phase~3: Analysis, Stretch Goals, and Dissemination (Months 16--24)}

\begin{itemize}
    \item \textbf{Data analysis:} Compare BO-optimized designs against naive
        baselines; quantify improvement in energy absorption.
    \item \textbf{Stretch goals:} If time permits, compute specific energy
        absorption (\si{J/kg}), effective spring constant normalized by weight
        ($k_{\text{eff}} / W$) and volume ($k_{\text{eff}} / V$).
    \item \textbf{Dissemination:} Prepare results for conference presentation
        (e.g., ASME IMECE, SES) and/or a journal publication.  Students will
        lead the writing and presentation.
\end{itemize}

% =============================================================================
\section{Mentoring Plan}
\label{sec:mentoring}
% =============================================================================

\textbf{The mentoring of undergraduate researchers is the central objective of
this proposal.}  The project is designed so that each student develops a broad
skill set while contributing to a cohesive research effort.

\subsection{Student Recruitment and Composition}

We plan to recruit 2--3 undergraduate students, ideally from courses in
mechanics of materials, machine design, kinematics, or compliant mechanisms.
Recruitment will include brief presentations in relevant classes to attract
students with appropriate background and genuine research interest.  We will
actively seek to recruit students from underrepresented groups in engineering.

\subsection{Mentoring Structure}

\begin{itemize}
    \item \textbf{Weekly meetings:} Individual 30-minute check-ins with each
        student plus a weekly 1-hour group meeting for progress updates, paper
        discussions, and skill-building workshops.
    \item \textbf{Graduate student co-mentor:} A portion of the budget will
        support a graduate student who provides day-to-day guidance, code
        reviews, and hands-on lab training, modeling effective mentoring
        practices.
    \item \textbf{Progressive responsibility:} Students will begin with
        structured tasks (running simulations with provided scripts, assembling
        structures from detailed CAD plans) and progress to independent design
        decisions and data interpretation.
    \item \textbf{Peer learning:} Students will present to each other weekly,
        developing scientific communication skills in a low-stakes environment
        before presenting at group meetings or conferences.
\end{itemize}

\subsection{Skills Development}

Each student will gain experience in multiple areas:

\begin{table}[h]
\centering
\begin{tabular}{@{}ll@{}}
\toprule
\textbf{Skill Area} & \textbf{Activities} \\
\midrule
Experimental design  & Planning test matrices, background measurements, controls \\
Data acquisition     & Force sensor operation, data logging, signal processing \\
Simulation           & FEA modeling, Python scripting, parameter sweeps \\
Statistical modeling & Gaussian processes, Bayesian optimization, uncertainty \\
CAD and fabrication  & SolidWorks/Fusion 360, 3D printing, assembly \\
Scientific communication & Lab notebooks, group presentations, conference talks \\
\bottomrule
\end{tabular}
\caption{Skill areas and corresponding student activities.}
\label{tab:skills}
\end{table}

\subsection{Assessment of Mentoring Effectiveness}

We will track mentoring outcomes through:
\begin{itemize}
    \item Structured mid-semester and end-of-semester self-assessments where
        students reflect on skills gained and confidence levels.
    \item Faculty evaluation of student independence, scientific reasoning, and
        communication growth.
    \item Tracking of tangible outcomes: conference abstracts submitted, poster
        presentations, and any manuscripts in preparation.
    \item Graduate school and career outcomes for participating students
        (tracked post-project).
\end{itemize}

% =============================================================================
\section{Timeline}
% =============================================================================

\begin{table}[h]
\centering
\small
\begin{tabular}{@{}lcccc@{}}
\toprule
\textbf{Task} & \textbf{Y1 Fall} & \textbf{Y1 Winter} & \textbf{Y2 Fall} & \textbf{Y2 Winter} \\
\midrule
Student recruitment \& onboarding    & \(\bullet\) &           &           &           \\
Literature review                     & \(\bullet\) & \(\bullet\) &           &           \\
Tensegrity parameterization           & \(\bullet\) & \(\bullet\) &           &           \\
Simulation model development          &             & \(\bullet\) & \(\bullet\) &           \\
Bayesian optimization implementation  &             & \(\bullet\) & \(\bullet\) &           \\
CAD design \& fabrication             &             &             & \(\bullet\) & \(\bullet\) \\
Impact testing \& data collection     &             &             & \(\bullet\) & \(\bullet\) \\
Simulation--experiment comparison     &             &             &             & \(\bullet\) \\
Stretch goal analysis                 &             &             &             & \(\bullet\) \\
Conference presentation / manuscript  &             &             &             & \(\bullet\) \\
\bottomrule
\end{tabular}
\caption{Project timeline across four semesters.}
\label{tab:timeline}
\end{table}

% =============================================================================
\section{Budget}
% =============================================================================

% sections/budget.tex
% Budget justification for BYU Mentored Research Grant

Total requested: \textbf{\$25,000} over 2 years.  Per grant guidelines, no
funds are allocated to equipment, tuition, or faculty wages.

\begin{center}
{\small
\begin{tabular}{@{}llrr@{}}
\toprule
\textbf{Category} & \textbf{Description} & \textbf{Year 1} & \textbf{Year 2} \\
\midrule
Undergraduate wages
    & 2 students, part-time + summer full-time
    & \$9,000 & \$9,000 \\
Graduate student wages
    & Partial RA for grad co-mentor
    & \$1,250 & \$1,250 \\
Supplies
    & Filament, load cells, DAQ, compute
    & \$2,000 & \$500 \\
Travel
    & Conference registration \& travel
    & \$0 & \$1,500 \\
Other
    & Poster printing, consumables
    & \$250 & \$250 \\
\midrule
\textbf{Annual Total} & & \textbf{\$12,500} & \textbf{\$12,500} \\
\bottomrule
\end{tabular}
}
\end{center}

\textbf{Undergraduate wages (\$18,000):} Supports 2 undergraduates working
part-time (\textasciitilde 10~hrs/week) during fall and winter semesters and
full-time (\textasciitilde 40~hrs/week) during summer terms for simulation,
3D printing, and testing.
%
\textbf{Graduate student wages (\$2,500):} Prof.~Baird's master's student
serves as day-to-day co-mentor.
%
\textbf{Supplies (\$2,500):} PLA/TPU filament, load cells, DAQ module,
fasteners, and cloud compute credits (Year~1 heavy).
%
\textbf{Travel (\$1,500):} Student presentations at ASME IDETC or UCUR
(Year~2).
%
\textbf{Other (\$500):} Poster printing, consumables, contingency.


% =============================================================================
\section{Broader Impacts}
% =============================================================================

\begin{itemize}
    \item \textbf{Undergraduate development:} This project will directly mentor
        2--3 undergraduates through a complete research arc, from literature
        review to conference dissemination, preparing them for graduate study or
        research-oriented careers.

    \item \textbf{Interdisciplinary training:} Students will work at the
        intersection of structural mechanics, computational modeling, and
        data-driven optimization---a combination increasingly valued in both
        academia and industry.

    \item \textbf{Open science:} Simulation code, optimization scripts, and
        experimental data will be made publicly available via GitHub, enabling
        reproducibility and extending impact beyond BYU.

    \item \textbf{Course integration:} Insights and demonstrations from this
        project may be incorporated into mechanics of materials or machine
        design courses, enriching the curriculum for future students.
\end{itemize}

% =============================================================================
\section{Qualifications}
% =============================================================================

\todo{Add PI qualifications, relevant publications, and prior mentoring
experience.  Emphasize past success mentoring undergraduates and any relevant
expertise in mechanics, optimization, or computational methods.}

% =============================================================================
\bibliographystyle{plainnat}
\bibliography{references}

\end{document}
